\headline{{\textbf{\Large\sffamily Seasonal Care Tips:}}\\
          {\textbf{\large\sffamily Mid\--December to Mid\--January}}}
\label{2025-12-care}
\vfill

\noindent
As we navigate the heart of the winter solstice, our collections enter their most vulnerable phase of dormancy.
The mild, damp conditions that characterised the first half of December are forecast to give way to a far more aggressive weather pattern as we transition into the New Year.
We are currently facing a ``tale of two winters'': a soggy, temperate conclusion to 2025, followed by a predicted sharp ``snow bomb'' and deep arctic airflow in the first week of January.
This shift from saturated roots to snap\--freezing conditions is a significant physiological stressor for bonsai.
Our primary focus for the coming month is managing the transition from moisture control to extreme thermal protection, ensuring that the structural integrity of both the trees and their ceramic containers remains uncompromised through the coming deep freeze.

\medskip
\topic{The Transition from Wet to Frozen}
\noindent
The immediate priority is the management of the current persistent rainfall.
With the ground and pots already heavily saturated, the forecasted drop in temperature in early January poses a dual threat: root rot from standing water and ``pot\--cracking'' frost heave.
If your trees are still exposed, now is the time to move them.
The transition from a mild 9°C to a sub\--zero ``snow bomb'' can happen in hours.
Ensure that all deciduous and hardy evergreen specimens are no longer sitting in saucers or flat on benches.
Use \textbf{pot feet or wooden slats} to provide an air gap.
This gap is essential; it prevents the pot from ``gluing'' itself to the bench via ice, which can lead to drainage holes being plugged by frozen water, effectively drowning the roots in an icy slush when the sun briefly hits the canopy.

\medskip
\topic{Deep Frost and Snow Protection}
\noindent
The first two weeks of January are predicted to bring the lowest temperatures of the season, with overnight lows potentially reaching −2°C or lower.
For our local microclimate, this often manifests as a ``black frost.'' \textbf{Hardy species} like Maples, Elms, and Larches must have their root balls protected.
If you cannot move them into a cold frame, wrap the pots in \textbf{multiple layers of horticultural fleece or bubble wrap}, ensuring the insulation extends over the rim of the pot to protect the nebari.
For those with \textbf{Satsuki Azaleas or Mediterranean species}, these must be moved into an unheated shed or greenhouse immediately.
While snow itself can act as a natural insulator, the weight of a heavy ``snow bomb'' can easily snap the brittle, dormant branches of refined specimens.
Use a soft brush to clear heavy accumulations from the canopy of wired trees to prevent structural failure.

\medskip
\topic{Watering in the Dead of Winter}
\noindent
Watering during this period requires a counter\--intuitive approach.
During the damp end of December, you will likely not need to water at all.
However, as the cold snap arrives in January, the air often becomes much drier. \textbf{Frozen soil creates physiological drought}; the tree cannot ``drink'' frozen water, but wind continues to wick moisture from the branches and buds.
Check your trees during the midday thaw.
If the soil surface feels dry and the frost has temporarily melted, provide a small amount of water.
Crucially, use water that has reached \textbf{ambient temperature}---never use water straight from an outdoor tap that may be near\--freezing, as the thermal shock to the roots can be fatal.
If a tree is bone\--dry but the root ball is a solid block of ice, move it to a slightly warmer spot (like a garage) to thaw slowly before attempting to hydrate.

\medskip
\topic{Pest Vigilance and Winter Hygiene}
\noindent
The damp conditions leading up to the New Year are a catalyst for fungal pathogens.
While the trees are bare, carry out a ``winter wash'' of the benches.
Fungal spores and \textbf{overwintering pests} like scale insect and aphid eggs tuck themselves into the crevices of rough bark or under the rims of pots.
A soft\--bristled toothbrush can be used to gently clean the trunks of older specimens.
Pay close attention to the soil surface; remove any moss that has become overly thick or ``mucky,'' as this can trap too much moisture against the trunk base, leading to bark rot.
If you notice any \textbf{liverwort} appearing in the damp conditions, remove it now before it releases spores in the spring.
Maintaining a clinical level of cleanliness now will significantly reduce your workload when the larvae emerge in March.

\medskip
\topic{The Art of Winter Observation}
\noindent
With active growth at a standstill, use this month for ``visual pruning.''
This is the best time to evaluate the \textbf{ramification and silhouette} of your deciduous trees without the distraction of foliage.
Take photographs of your trees from all four sides.
These images are invaluable for planning spring repotting angles and structural wire placement.
If you see any wire that was applied in autumn beginning to look tight---highly unlikely now, but possible on vigorous species---remove it.
Otherwise, resist the temptation to prune.
Every cut made in mid\--winter is a wound that will not begin to heal for another eight weeks, leaving an entry point for silver leaf fungus or dieback during the January freezes.

\medskip
\topic{Preparing for the Spring Rush}
\noindent
As we pass the solstice, the days will slowly begin to lengthen.
While it feels like deep winter, the ``sap rise'' is only two months away for some species.
Use the quiet, cold days of early January to \textbf{sift and mix your substrates}.
Preparing your Akadama, Pumice, and Kiryu mixes now ensures they are dry and ready for the first repotting window in late February.
Check your wire stocks and order any specialty tools you might need.
This period of preparation is the foundation of a successful growing season.
By ensuring your trees are protected from the coming January ``snow bomb'' and your workshop is organised, you can enjoy the quiet beauty of the winter silhouette with the peace of mind that your collection is secure.
\closearticle
\columnbreak
